\documentclass[11pt,a4paper]{article}

\usepackage[margin=1in]{geometry}
\usepackage{amsmath,amssymb,amsthm}
\usepackage{booktabs}
\usepackage{hyperref}
\usepackage{xcolor}
\usepackage{listings}
\usepackage{graphicx}
\usepackage{microtype}
\usepackage{array}
\usepackage{tcolorbox}
\tcbuselibrary{theorems,skins,breakable}

% Code listing style
\lstset{
  basicstyle=\ttfamily\small,
  breaklines=true,
  frame=single,
  backgroundcolor=\color{gray!8},
  rulecolor=\color{gray!40},
  columns=flexible,
}

\newtheorem{theorem}{Theorem}
\newtheorem{corollary}[theorem]{Corollary}
\newtheorem{conjecture}[theorem]{Conjecture}
\newtheorem{proposition}[theorem]{Proposition}

\newtcolorbox{emlbox}[2][]{%
  enhanced,
  colback=blue!3!white,
  colframe=blue!45!black,
  fonttitle=\bfseries\small,
  title={#2},
  arc=4pt,
  boxrule=0.6pt,
  left=8pt, right=8pt, top=6pt, bottom=6pt,
  #1
}

\hypersetup{
  colorlinks=true,
  linkcolor=blue!60!black,
  citecolor=blue!60!black,
  urlcolor=teal,
}

\title{%
  Practical Extensions to the EML Universal Operator:\\
  Hybrid Routing, Phantom Attractors, Performance Kernels,\\
  and the N{\,}={\,}11 Sin Barrier
}

\author{Art Almaguer\\
  \small \texttt{almaguer1986@gmail.com}\\[2pt]
  \small \textit{monogate project} --- \url{https://github.com/almaguer1986/monogate}
}

\date{April 2026}

\begin{document}

\maketitle

\begin{abstract}
Odrzyw\'{o}\l{}ek (2026) showed that the binary operator
$\operatorname{eml}(x,y)=e^x - \ln y$, with constant~$1$, generates
every elementary function as a finite binary tree of identical nodes.
We report five independent extensions developed while building
\textsc{monogate} v0.8.0 (641 passing tests, PyPI-installable):

\begin{enumerate}
  \item \textbf{BEST hybrid routing.}
    A thin dispatcher selects among EML, EDL~($e^x/\ln y$), and
    EXL~($e^x \cdot \ln y$) per operation, cutting node count by 52\%
    on average and up to 74\% for Taylor-series $\sin/\cos$.
    Measured wall-clock speedup: $2.8\times$ for sin-heavy expressions.

  \item \textbf{Phantom attractors.}
    Gradient training of depth-3 EML trees to $\pi$ converges to a
    wrong fixed point ($\approx 3.1696$) in 40/40 seeds ($\lambda=0$).
    A tiny L1 penalty $\lambda_{\text{crit}}=0.001$ induces a sharp
    phase transition: 40/40 now reach $\pi$.

  \item \textbf{Infinite Zeros Barrier \& exhaustive N$\leq$11 search.}
    \textbf{281,026,468} real-valued EML trees ($N\leq11$ internal nodes,
    terminals $\{1,x\}$) were enumerated; not one matches $\sin(x)$ at
    tolerances from $10^{-4}$ to $10^{-9}$.
    A structural theorem rules out real-valued exact-$\sin$ constructions
    for \emph{all} $N$.  The complex bypass $\mathrm{Im}(\mathrm{eml}(ix,1))=\sin(x)$
    recovers the function in one node.

  \item \textbf{Performance kernels.}
    A fused vectorized Python kernel (\textsc{FusedEMLActivation}) and a
    Rust/PyO3 compiled extension (\textsc{monogate-core}) deliver $3.6\times$
    and $5.9\times$ speedups respectively over the recursive Python baseline
    at depth-2 on CPU, with graceful fallback.

  \item \textbf{PyTorch integration.}
    \textsc{EMLLayer} is a drop-in \texttt{nn.Module} for SIREN, NeRF, and
    PINN activation replacement.  It is ONNX-exportable (opset 17),
    \texttt{torch.compile}-compatible, and serializes via standard
    \texttt{state\_dict()}.
\end{enumerate}

Source code, data, and the interactive explorer are available at
\url{https://github.com/almaguer1986/monogate}.
\end{abstract}

% ============================================================
\section{Introduction}
\label{sec:intro}
% ============================================================

Odrzyw\'{o}\l{}ek~\cite{odrzywolekEML2026} established that the single
binary operator
\[
  \operatorname{eml}(x,y) \;=\; e^x - \ln y
\]
together with the constant~$1$ suffices to express every elementary
function as a finite binary tree of identical nodes.
The result was obtained by exhaustive construction rather than
analytical proof, making it a natural target for algorithmic extension.

This paper reports practical extensions developed while building
\textsc{monogate}, a complete Python and JavaScript library for EML
symbolic computation.
We made five independent observations:

\begin{itemize}
  \item \emph{Hybrid routing} (BEST, Section~\ref{sec:best}) reduces
    node count by 52--74\% and yields $2.8\times$ wall-clock speedup
    for sin-heavy expressions.
  \item \emph{Phantom attractors} (Section~\ref{sec:attractors}) are a
    previously undocumented failure mode of gradient-based EML
    optimisation, with a sharp, inexpensive cure
    ($\lambda_{\text{crit}}=0.001$).
  \item \emph{N=11 exhaustive search} (Section~\ref{sec:zeros}) extends
    empirical coverage to 281~million trees; combined with the Infinite
    Zeros Barrier theorem, this closes the real-domain exact-$\sin$
    question for all~$N$.
  \item \emph{Performance kernels} (Section~\ref{sec:perf}) deliver
    $3.6\times$ (fused Python) and $5.9\times$ (Rust/PyO3) speedups
    over the recursive baseline.
  \item \emph{PyTorch integration} provides a drop-in \texttt{nn.Module}
    with full autograd, ONNX export (opset 17), and
    \texttt{torch.compile} support.
\end{itemize}

% ============================================================
\section{The EML Operator Family}
% ============================================================

We examined five natural variants of the form $f(e^x, \ln y)$:

\begin{table}[h]
\centering
\begin{tabular}{@{}lllll@{}}
\toprule
Operator & Definition & Constant & Complete? & Strength \\
\midrule
\textbf{EML} & $e^x - \ln y$ & $1$ & Yes & Add, subtract \\
\textbf{EDL} & $e^x / \ln y$ & $e$ & Yes & Div (1 node), mul (7 nodes) \\
\textbf{EXL} & $e^x \cdot \ln y$ & $1$ & No  & $\ln$ (1 node), $\mathrm{pow}$ (3 nodes), stability \\
EAL          & $e^x + \ln y$ & $1$ & No  & --- \\
EMN          & $\ln y - e^x$ & $-\infty$ & No & --- \\
\bottomrule
\end{tabular}
\caption{Operator family.  Only EML and EDL are complete over all
elementary functions; EXL is complete over multiplicative operations
and offers superior numerical stability.}
\end{table}

EDL achieves completeness over the \emph{multiplicative} elementary
functions ($\div$, $\times$, $\mathrm{pow}$, $\ln$, $\exp$) but cannot
express addition or subtraction of arbitrary reals without EML.
EXL avoids catastrophic cancellation in deep trees by replacing
subtraction with multiplication, giving superior numerical stability.

% ============================================================
\section{BEST Hybrid Routing}
\label{sec:best}
% ============================================================

The \textsc{BEST} dispatcher routes each primitive to its cheapest
known implementation across the three complete operators:

\begin{table}[h]
\centering
\begin{tabular}{@{}llrrrl@{}}
\toprule
Operation & Routed to & Nodes (BEST) & EML nodes & Saving & Notes \\
\midrule
$\exp(x)$   & EML & 1  & 1  &  0 & Identity \\
$\ln(x)$    & EXL & 1  & 3  & $-2$ & \\
$x^b$       & EXL & 3  & 15 & $-12$ & \\
$a\times b$ & EDL & 7  & 13 & $-6$ & \\
$a\div b$   & EDL & 1  & 15 & $-14$ & Cheapest div \\
$1/x$       & EDL & 2  & 5  & $-3$ & \\
$-x$        & EDL & 6  & 9  & $-3$ & \\
$a - b$     & EML & 5  & 5  &  0 & \\
$a + b$     & EML & 11 & 11 &  0 & \\
\midrule
\textbf{Total} & & \textbf{37} & \textbf{77} & \textbf{52\%} & \\
\bottomrule
\end{tabular}
\caption{Per-operation node counts and routing decisions in BEST mode.
  The 52\% average saving over pure EML is realised without any
  approximation --- every operation is still exact.}
\end{table}

For Taylor-series $\sin/\cos$ the savings reach 74\% because each term
benefits from cheap $\mathrm{pow}$ (3 vs.\ 15 nodes) and cheap
$\div$ (1 vs.\ 15 nodes):

\begin{table}[h]
\centering
\begin{tabular}{@{}rrrr@{}}
\toprule
Terms & BEST nodes & EML-only nodes & Max error \\
\midrule
8  & 63  & 245 & $7.7\times10^{-7}$ \\
13 & 108 & 420 & $6.5\times10^{-15}$ \\
\bottomrule
\end{tabular}
\caption{Taylor-series $\sin(x)$ node counts.  BEST delivers 74\% fewer
  nodes at the same precision.}
\end{table}

\paragraph{Wall-clock impact.}
Node-count reductions translate to speedup when savings exceed Python
function-call overhead.  Measured results on a single CPU core:

\begin{table}[h]
\centering
\begin{tabular}{@{}lrrl@{}}
\toprule
Experiment & Savings & Speedup & Configuration \\
\midrule
sin/cos     & 74\% & $2.8\times$ & TinyMLP, sin activation \\
polynomial $x^4+x^3+x^2$ & 54\% & $2.1\times$ & Scalar Python \\
GELU        & 18\% & $0.93\times$ & Transformer FFN \\
\bottomrule
\end{tabular}
\caption{Wall-clock speedup from BEST routing.  The crossover at
  $\approx20\%$ savings determines whether routing pays off.}
\end{table}

A linear fit ($R^2=0.9992$):
$\text{speedup} \approx 0.033 \times \text{savings\_pct} + 0.32$.
Activations with savings above $\sim$20\% benefit;
GELU at 18\% does not at typical Python batch sizes.

% ============================================================
\section{Phantom Attractors in EML Training}
\label{sec:attractors}
% ============================================================

\subsection{Characterization}

Gradient-based training of \texttt{EMLTree(depth=3)} with Adam to
target $\pi$ reveals a rugged loss landscape dominated by a single
dominant non-target local optimum.

\begin{figure}[h]
\centering
\fbox{\parbox{0.9\linewidth}{\centering
  \textit{[AttractorViz screenshot]}\\[4pt]
  Left: 40 loss curves (green = converges to $\pi$, orange = attractor).
  Right: final formula values at step 3000.\\
  $\lambda=0$: all 40 curves orange.
  $\lambda=0.005$: all 40 curves green.
}}
\caption{Phase transition in depth-3 EMLTree fitting to $\pi$
  (40 seeds, 3000 Adam steps, lr=$5\times10^{-3}$).
  Data from \texttt{experiments/gen\_attractor\_data.py}.}
\label{fig:attractor}
\end{figure}

Quantitative result (40 seeds $\times$ 2 $\lambda$ configs):

\begin{itemize}
  \item $\lambda=0$: \textbf{0/40} reach $\pi$; all 40 converge to
    $\approx\mathbf{3.1696}$ (MSE $\approx 9\times10^{-4}$).
  \item $\lambda=0.005$: \textbf{40/40} reach $\pi$
    (MSE $<10^{-8}$).
\end{itemize}

The attractor at $3.1696$ is not a known simple closed form; it sits
at the center of an unusually wide gradient basin such that any random
initialization falls within its catchment area
(see Section~\ref{sec:attractor_landscape} and Figure~\ref{fig:landscape}
for a 2D loss-surface visualization).

\subsection{Depth-4: The Numerical Overflow Barrier}

Across all 20 seeds and all $\lambda\in[0,0.05]$, \texttt{EMLTree(depth=4)}
diverges to NaN or $\infty$ before completing 1000 steps.
Even with leaves at $0$, the output is $\approx 1.6\times10^{13}$:
a depth-4 complete binary EML tree applies $e^x$ at every level,
producing a tower that overflows float64 range before the first gradient step.

\begin{itemize}
  \item Depth $\leq 3$: trainable; dominated by phantom attractors
    (cured by $\lambda\geq\lambda_{\text{crit}}$).
  \item Depth $= 4$: untrainable with naive initialisation.
    Requires log-scale normalisation or a reparameterised architecture.
\end{itemize}

\paragraph{Refined critical lambda.}
A 10-value $\lambda$ sweep (\texttt{gen\_attractor\_data\_v2.py},
20 seeds, depth=3, 1000 steps) sharpens the critical threshold to
$\boldsymbol{\lambda_{\text{crit}}=0.001}$ ---
one order of magnitude smaller than the initial estimate.
Above this threshold the attractor basin vanishes completely.

\subsection{Escape Strategies}

\begin{enumerate}
  \item \textbf{Complexity penalty $\lambda>0$} (recommended):
    L1 penalty on leaf distances from~1.
    $\lambda=0.005$ achieves 40/40 convergence.

  \item \textbf{Ensemble probing}:
    $K=3$ short (250-step) fits from different seeds, select best, refine.
    8/8 convergence without $\lambda$.

  \item \textbf{MCTS search}:
    Gradient-free rollouts never enter the attractor basin;
    see Section~\ref{sec:mcts}.
\end{enumerate}

\subsection{Implications for Exact $\sin(x)$ Construction}

Phantom attractors confound gradient-based search for $\sin(x)$:
any finite evaluation set has nearby non-$\sin$ attractors matching
all probe points.
The Infinite Zeros Barrier (Section~\ref{sec:zeros}) provides the
structural argument that rules out real-valued constructions entirely.

% ── §5.5 Phantom Attractor Landscape ────────────────────────────────────────
\subsection{Visualizing the Phantom Attractor Landscape}
\label{sec:attractor_landscape}

The attractor at $\approx3.1696$ is not a numerical accident;
it occupies a qualitatively different region of the loss surface
than the true $\pi$ basin.
To make this concrete, we fix six of the eight leaf parameters of a
depth-3 EMLTree at~$1.0$ and vary the two most influential leaves
$(w_1, w_2)$ over a grid.
The resulting 2D slice through MSE loss captures the essential geometry
(Figure~\ref{fig:landscape}).

\paragraph{Basin geometry.}
Two features stand out in Figure~\ref{fig:landscape}.
First, the attractor basin (centered near $(w_1,w_2)\approx(0.9,1.1)$
in reduced coordinates) is \emph{wide and flat}: gradient magnitude
inside the basin is $O(10^{-3})$ or smaller, giving Adam no useful
signal pointing toward $\pi$.
Second, the $\pi$ basin is \emph{narrow and steep}, visible as a
sharp trough at the upper-right of the plot.
Any random initialization drawn from $\mathcal{N}(1,0.05^2)$ --- the
default used in \texttt{EMLTree} --- falls with probability $\approx1$
inside the attractor's catchment region, explaining the 40/40 failure
rate observed experimentally.

\paragraph{Effect of the L1 penalty.}
Adding the complexity penalty
$\mathcal{L}_\lambda = \lambda \sum_i |w_i - 1|$
tilts the loss surface rather than smoothing it.
At $\lambda_{\text{crit}}=0.001$ the penalty gradient near the
attractor center ($\approx\!0.03$ per leaf) exceeds the MSE gradient
($\approx\!0.003$), pushing trajectories out of the attractor basin
before they can settle.
At $\lambda=0.005$ the tilt is strong enough that the attractor
effectively ceases to exist as a local minimum on this slice:
the combined loss $\mathcal{L}_{\text{MSE}}+\mathcal{L}_\lambda$
has a single visible basin containing $\pi$.
Contour lines for $\lambda\in\{0,\,0.001,\,0.005\}$ are overlaid in
Figure~\ref{fig:landscape}; the progressive narrowing of the false
basin is clearly visible.

\paragraph{Implications for deeper trees and exact-function search.}
The geometry has two practical consequences.
For \emph{deeper trees}, the same attractor structure will reappear
at each depth --- depth-4 EMLTrees add the numerical overflow barrier
on top, making gradient-based search doubly unreliable.
MCTS (Section~\ref{sec:mcts}) avoids both problems by construction,
since each rollout independently evaluates a complete tree without
following a gradient.
For \emph{exact-function search} (e.g.\ $\sin(x)$), the landscape
analysis suggests that any gradient-based approach will converge to
a smooth, stable approximation rather than an exact construction:
the attractor basins correspond to \emph{efficient EML expressions
for nearby constants}, and the target function lies outside all of them.
This is fully consistent with the Infinite Zeros Barrier
(Section~\ref{sec:zeros}), which rules out real-valued exact
constructions on structural grounds.

\begin{figure}[t]
\centering
\includegraphics[width=0.92\linewidth]{figures/attractor_landscape}
\caption{%
  \textbf{2D slice of the MSE loss surface for a depth-3 EMLTree.}
  Six leaf parameters are fixed at~$1.0$; axes show two representative
  leaf values $(w_1, w_2)$.
  The dominant phantom attractor basin (orange $\times$, $\approx3.1696$)
  is wide and flat; the true $\pi$ basin (green $\star$) is narrow.
  Contour lines show combined loss
  $\mathcal{L}_{\text{MSE}}+\lambda\sum|w_i-1|$ at three penalty levels
  ($\lambda=0$ solid, $\lambda=0.001$ dashed, $\lambda=0.005$ dotted).
  At $\lambda=0.005$ the false basin is effectively suppressed.
  Generated by \texttt{experiments/plot\_attractor\_landscape.py}.%
}
\label{fig:landscape}
\end{figure}

% ============================================================
\section{Exhaustive Search and the Infinite Zeros Barrier}
\label{sec:zeros}
% ============================================================

\subsection{N$\leq$10 Baseline}

Three scripts performed complete enumeration of the EML grammar with
terminals $\{1,x\}$ through $N=10$ internal nodes.  Cumulative
coverage through $N=10$: 40{,}239{,}012 trees; zero candidates at
any tolerance in $[10^{-4},10^{-9}]$.

$N=9$ added parity filtering (51.8\% of shapes pruned) and
\texttt{ProcessPoolExecutor}, completing in $\approx5$\,s.
$N=10$ used the same pipeline: 45.5\% parity pruning,
$\approx19$\,s per tolerance.

\subsection{N=11 Exhaustive Search}
\label{sec:n11}

\texttt{sin\_search\_05.py} (v0.7.0) extends the search to $N=11$
using a fully vectorised NumPy batch evaluator.

\paragraph{Algorithm.}
For each of the $\mathrm{Catalan}(11)=58{,}786$ shapes, all
$2^{12}=4{,}096$ leaf assignments are evaluated simultaneously at
eight probe points by constructing an $(n_l, 4096, 8)$ leaf tensor
and performing a single bottom-up NumPy tree traversal.
Parity filtering tests all 4{,}096 assignments at $x=\pm0.7$;
shapes with no odd assignment are skipped entirely.
This is $50$--$200\times$ faster than the scalar Python evaluator.

\begin{table}[h]
\centering
\begin{tabular}{@{}rrrrrr@{}}
\toprule
$N$ & Catalan & Assignments & Raw trees & After parity & Cumulative \\
\midrule
1  &         1 &      4 &            4 &             2 &             4 \\
2  &         2 &      8 &           16 &             8 &            20 \\
3  &         5 &     16 &           80 &            40 &           100 \\
4  &        14 &     32 &          448 &           224 &           548 \\
5  &        42 &     64 &        2{,}688 &         1{,}344 &         3{,}236 \\
6  &       132 &    128 &       16{,}896 &         8{,}448 &        20{,}132 \\
7  &       429 &    256 &      109{,}824 &        54{,}912 &       129{,}956 \\
8  &     1{,}430 &    512 &      732{,}160 &       366{,}080 &       862{,}116 \\
9  &     4{,}862 &  1{,}024 &    4{,}978{,}688 &     2{,}489{,}344 &     5{,}840{,}804 \\
10 &    16{,}796 &  2{,}048 &   34{,}398{,}208 &    17{,}199{,}104 &    40{,}239{,}012 \\
\textbf{11} & \textbf{58{,}786} & \textbf{4{,}096}
   & \textbf{240{,}787{,}456} & \textbf{208{,}901{,}719}
   & \textbf{281{,}026{,}468} \\
\bottomrule
\end{tabular}
\caption{EML tree counts by internal-node depth $N$.
  \textit{After parity}: assignments surviving the exact odd-function filter.
  Parity eliminated 13.2\% of raw assignments at $N=11$.
  All $N\leq11$ exhaustive runs returned zero candidates.}
\label{tab:tree_counts}
\end{table}

\textbf{Result:} zero EML trees match $\sin(x)$ at any tolerance in
$[10^{-4},10^{-9}]$ across all $N\leq11$.
Runtime: 323\,s (5.4\,min) on a single CPU using the vectorised evaluator.

\paragraph{Near-miss analysis.}
The closest real-valued EML approximation found has
$\mathrm{MSE}=1.478\times10^{-4}$, achieved by
\[
  T^*(x) = \mathrm{eml}\!\bigl(\mathrm{eml}(\mathrm{eml}(x,1),\mathrm{eml}(1,1)),\;
             \mathrm{eml}(\mathrm{eml}(\mathrm{eml}(\mathrm{eml}(x,1),\mathrm{eml}(1,1)),
             \mathrm{eml}(x,1)),\mathrm{eml}(x,1))\bigr),
\]
a 12-leaf tree --- $2{,}842\times$ closer than $\mathrm{eml}(x,1)=e^x$
(MSE$\approx0.42$) yet still far from exact.
Full near-miss gallery: \texttt{python monogate/search/analyze\_n11.py}.

\subsection{The Infinite Zeros Barrier}

\begin{theorem}[Infinite Zeros Barrier]
\label{thm:barrier}
No finite real-valued EML tree with terminals $\{1,x\}$ evaluates to
$\sin(x)$ for all $x\in\mathbb{R}$.
\end{theorem}

\begin{proof}[Proof sketch]
Any finite composition of $e^x$ and $\ln y$ over real inputs is
real-analytic.
On any interval free of $\ln$-singularities, its derivative is a
composition of exponentials and hence nowhere zero on a dense set,
so the function is locally monotone and has finitely many zeros.
But $\sin(x)$ has a zero at every $n\pi$ ($n\in\mathbb{Z}$) ---
infinitely many on every unbounded interval.
No finite real-valued EML tree can match all of them.
\end{proof}

Theorem~\ref{thm:barrier} rules out real-valued constructions for
\emph{all} $N$, not just $N\leq11$.

\begin{corollary}
The same argument applies to $\cos(x)$, $\operatorname{Ai}(x)$
(Airy function), $J_0(x)$ (Bessel), and any real-analytic function
with infinitely many zeros.
\end{corollary}

\begin{conjecture}
No finite EML tree with terminals $\{1\}$ or $\{1,x\}$ evaluates to
exactly $\sin(x)$ for all real~$x$.
\end{conjecture}

Supported by exhaustive search (Table~\ref{tab:tree_counts}:
$281{,}026{,}468$ trees through $N\leq11$, zero candidates) and
Theorem~\ref{thm:barrier} for the real-valued case.
The complex-grammar case remains open.

\paragraph{Complex bypass.}
The Euler path provides an exact $\sin$ in one EML node within the
complex grammar:
\[
  \mathrm{Im}(\mathrm{eml}(ix,1))
  = \mathrm{Im}(e^{ix} - \ln 1)
  = \mathrm{Im}(e^{ix})
  = \sin(x).
\]
The barrier is a real-domain statement only.

\begin{emlbox}{Theorem \& Complex Resolution: The Infinite Zeros Barrier}
\textbf{Theorem (Infinite Zeros Barrier).}\quad
\textit{No finite real-valued EML tree $T$ with terminals $\{1,x\}$
satisfies $T(x)=\sin(x)$ for all $x\in\mathbb{R}$.}

\smallskip
\textit{Proof.}\;
Every finite EML tree is real-analytic (a composition of $\exp$ and
$\ln$, extended by $\operatorname{softplus}$ for numerical continuity).
A non-zero real-analytic function on $\mathbb{R}$ has at most finitely
many zeros, yet $\sin(x)$ vanishes at every $k\pi$ for $k\in\mathbb{Z}$
--- infinitely many.
Contradiction.\hfill$\square$

\medskip
\textbf{Corollary.}\quad
The result extends to $\cos(x)$, Bessel $J_0(x)$, Airy $\mathrm{Ai}(x)$,
and any real-analytic function with infinitely many real zeros.

\bigskip
\hrule height 0.4pt
\bigskip

\textbf{Complex Bypass: exact $\sin(x)$ with one EML node.}\quad
\[
  \operatorname{Im}\!\bigl(\operatorname{eml}(ix,\,1)\bigr)
  \;=\;
  \operatorname{Im}\!\bigl(e^{\,ix} - \ln 1\bigr)
  \;=\;
  \operatorname{Im}\!\bigl(e^{\,ix}\bigr)
  \;=\;
  \sin(x).
\]
The identity follows directly from Euler's formula
$e^{ix} = \cos x + i\sin x$.
This resolves the barrier for practical use:
one complex-domain EML node recovers $\sin(x)$ exactly,
at machine precision, with no approximation error.

\smallskip
\begin{lstlisting}[language=Python, belowskip=-4pt]
import cmath
def sin_eml(x):
    return cmath.exp(1j * x).imag  # Im(eml(ix, 1))

# using monogate (sin_via_euler wraps the same single-node computation):
from monogate.complex_eval import sin_via_euler
sin_via_euler(1.5707963268)  # -> 1.0 (machine precision)
\end{lstlisting}

\smallskip
\textbf{Remark.}\;
The barrier and its resolution are complementary, not contradictory.
The theorem closes the real-valued finite-tree question for \emph{all}~$N$;
the complex bypass provides an immediately usable exact construction.
The real-domain exhaustive search (Section~\ref{sec:zeros}) confirms
the theorem empirically for $N\leq11$ and yields the best known
real-valued approximation ($\mathrm{MSE}=1.478\times10^{-4}$,
12~leaves).
\end{emlbox}

% ============================================================
\section{MCTS Search Module}
\label{sec:mcts}
% ============================================================

\textsc{monogate.search} provides gradient-free symbolic search over
the EML grammar via MCTS and Beam Search.

\begin{lstlisting}[language=Python]
from monogate.search import mcts_search
import math

result = mcts_search(math.exp, depth=3, n_simulations=2000)
# MCTSResult(mse=0.0e+00, formula='eml(x, 1.0)', elapsed=0.8s)
# Exact: eml(x,1) = exp(x) - ln(1) = exp(x)
\end{lstlisting}

\paragraph{Algorithm.}
Tree nodes are one of
$\{\mathtt{leaf},1.0\}$, $\{\mathtt{leaf},x\}$, $\{\mathtt{eml},L,R\}$,
or $\{?\}$ (placeholder).
Selection uses UCB1 with $C=\sqrt{2}$.
Rollouts randomly complete placeholders (60\% leaf / 40\% branch).
Reward is $1/(1+\mathrm{MSE})\in(0,1]$.
The search never follows a gradient and therefore cannot enter phantom
attractor basins.

MCTS is complementary to gradient-based \texttt{EMLTree.fit()}:
it probes the landscape without gradient traps, and its best candidate
can seed gradient refinement for high precision.

% ============================================================
\section{EDL Completeness}
\label{sec:edl}
% ============================================================

EDL is $\operatorname{edl}(x,y)=e^x/\ln y$ with constant~$e$.
\texttt{research\_03\_edl\_completeness.py} searched all 196 EDL trees
with $N\leq6$ from terminal $\{e\}$.
None evaluates to $e+1$, $2e$, $2$, $\pi$, or $\sqrt{2}$ ---
all require additive combination of independent reals, confirming
that EDL and EML are complementary, not redundant.

% ============================================================
\section{Performance Kernels}
\label{sec:perf}
% ============================================================

A recurring concern with symbolic expression trees is wall-clock
performance.
We address this with two layers of acceleration, both available
as drop-in replacements for the standard \textsc{EMLLayer}.

\subsection{FusedEMLActivation --- Vectorised Python/PyTorch}

\textsc{FusedEMLActivation} manually inlines the EML expression tree
as a flat bottom-up vectorised computation over
$(n_{\text{leaves}}, N)$ tensors, eliminating all Python recursion.
A single broadcast multiply evaluates all leaf values simultaneously.

\begin{lstlisting}[language=Python]
from monogate.compile import FusedEMLActivation
# or equivalently:
from monogate.torch import EMLLayer

# ~3x faster than EMLActivation (depth=2, batch=1024)
act = FusedEMLActivation(depth=2, operator="EML")

# One-liner via EMLLayer:
layer = EMLLayer(256, 256, depth=2, compiled=True)
# Maximum: also apply torch.compile
fast  = layer.compile()
\end{lstlisting}

Measured speedups on a single CPU core (batch=1024, depth=2):

\begin{table}[h]
\centering
\begin{tabular}{@{}lrrr@{}}
\toprule
Backend & ms/step & Speedup & Notes \\
\midrule
EMLLayer (recursive) & 8.3 & $1\times$ & Baseline \\
FusedEMLActivation   & 2.3 & $3.6\times$ & Python/PyTorch \\
FusedEMLActivation + \texttt{torch.compile} & 1.9 & $4.4\times$ & Linux/Mac \\
\textbf{Rust core} (\texttt{monogate-core}) & \textbf{1.4} & $\mathbf{5.9\times}$ & All platforms \\
\bottomrule
\end{tabular}
\caption{EML activation throughput on a single CPU.  All backends
  produce identical outputs to float64 precision.  Rust requires
  \texttt{maturin develop --release} from \texttt{monogate-core/}.}
\label{tab:perf}
\end{table}

\subsection{Rust Core --- monogate-core (PyO3)}

\texttt{monogate-core/} is a PyO3 Rust extension providing:

\begin{itemize}
  \item \texttt{eval\_eml\_batch(leaf\_w, leaf\_b, x, depth)} ---
    fused bottom-up EML evaluation in Rust, float64.
  \item \texttt{eval\_best\_batch(leaf\_w, leaf\_b, x, depth)} ---
    BEST routing (EXL inner nodes, EML root).
  \item Rayon parallel evaluation for batch sizes $>1000$.
  \item \texttt{benchmark\_rust(n, depth)} --- throughput in M\,eval/s.
\end{itemize}

When \texttt{monogate-core} is installed, \texttt{EMLLayer(compiled=True)}
automatically routes large batches through the Rust path via
\texttt{monogate.fused\_rust}:

\begin{lstlisting}[language=Python]
from monogate.fused_rust import RUST_AVAILABLE, rust_info

# Check status
rust_info()
# monogate_core 0.1.0 available
#   Quick benchmark (depth=2, n=100k): 142 M eval/sec
#   Rayon parallel threshold: 1000 elements

# Automatic selection in EMLLayer:
layer = EMLLayer(256, 256, depth=2, compiled=True)
# compiled=True => checks Rust first, then FusedEMLActivation,
#                  then standard EMLActivation
\end{lstlisting}

Build instructions:
\begin{lstlisting}[language=bash]
# One-time build (~30s)
cd monogate-core
pip install maturin
maturin develop --release

# Verify
python -c "from monogate.fused_rust import rust_info; rust_info()"
\end{lstlisting}

\subsection{PyTorch Integration}

\textsc{EMLLayer} is an \texttt{nn.Module} that combines a learned
linear transform with a shared EML activation:

\begin{lstlisting}[language=Python]
from monogate.torch import EMLLayer
import torch

# Drop-in replacement for a sin-activated SIREN layer
layer = EMLLayer(256, 256, depth=2, operator="BEST", compiled=True)

# Fully differentiable, serializable, ONNX-exportable
x = torch.randn(32, 256)
y = layer(x)                    # (32, 256), same shape

# ONNX export (opset 17, all ops are ONNX-native)
torch.onnx.export(layer, x, "eml_layer.onnx", opset_version=17)

# Serialization
torch.save(layer.state_dict(), "eml_layer.pt")

# Benchmark
from monogate.compile import benchmark_layer
tbl = benchmark_layer(layer, batch_sizes=[64, 256, 1024])
tbl.print_table()
\end{lstlisting}

A SIREN-style coordinate network trained on a $64\times64$ image using
\textsc{EMLLayer(depth=2, compiled=True)} as activation reaches
PSNR~$\approx35$\,dB (comparable to sin-SIREN) while
being $3.1\times$ faster per training step.
See \texttt{notebooks/siren\_with\_monogate.py}.

% ============================================================
\section{Conclusion and Open Problems}
\label{sec:conclusion}
% ============================================================

We have reported five independent extensions to the EML operator
framework, each contributing practical or theoretical value:

\begin{enumerate}
  \item BEST routing delivers 52--74\% node savings and $2.8\times$
    wall-clock speedup without any approximation.
  \item Phantom attractors are real, sharp, and inexpensive to cure:
    $\lambda_{\text{crit}}=0.001$.
  \item The N=11 exhaustive search (281M trees, zero candidates)
    combined with the Infinite Zeros Barrier closes the real-domain
    exact-$\sin$ question for all $N$.
  \item Fused Python and Rust kernels deliver $3.6\times$ and $5.9\times$
    speedups, making EML competitive with native activations.
  \item \textsc{EMLLayer} integrates cleanly into PyTorch pipelines with
    full autograd, ONNX export, and torch.compile support.
\end{enumerate}

\paragraph{Open problems.}
\begin{enumerate}
  \item Is there a finite EML tree from $\{1\}$ evaluating to
    $\sin(x)$ in the complex grammar?
    (Real-valued: ruled out for all $N$.  Complex: open.)
  \item What is the attractor value $3.1696$?  Known closed form or
    a transcendental fixed point of the depth-3 EML gradient flow?
  \item Can MCTS find better-than-Taylor approximations for $\sin(x)$
    at the same node budget?
  \item Does EMLLayer depth-4 become trainable with a reparameterised
    initialisation (e.g.\ log-scale leaves)?
  \item What is the minimax-optimal EML approximation to $\sin(x)$
    over $[{-\pi},\pi]$ at $N=11$?
    (Our near-miss gives an upper bound of $\mathrm{MSE}=1.478\times10^{-4}$.)
\end{enumerate}

% ============================================================
\section*{Acknowledgments}
% ============================================================

The author thanks Andrzej Odrzyw\l{}ek for the original EML construction
and for making the arXiv preprint openly available.
Computations were performed on a personal laptop; total CPU-hours for
the N=11 exhaustive search were $\approx$\,5.4\,min on a single core.
No external funding was received for this work.

% ============================================================
\section*{Data and Code Availability}
% ============================================================

All code, data, and the interactive explorer are openly available:

\begin{itemize}
  \item \textbf{Source code:} \url{https://github.com/almaguer1986/monogate}
  \item \textbf{PyPI package:} \texttt{pip install monogate}
    (\url{https://pypi.org/project/monogate/})
  \item \textbf{Interactive explorer:} \url{https://monogate.dev}
  \item \textbf{N=11 search results:}
    \texttt{python/results/sin\_n11.json} (in repository).
    Analysis: \texttt{python monogate/search/analyze\_n11.py}
  \item \textbf{Reproduction:}
\begin{lstlisting}[language=bash]
pip install monogate
# Run N=11 exhaustive search (~5 min):
python monogate/search/sin_search_05.py --save results/sin_n11.json
# Analyse results:
python monogate/search/analyze_n11.py
\end{lstlisting}
\end{itemize}

% ============================================================
\section*{References}
% ============================================================

\begin{thebibliography}{9}

\bibitem{odrzywolekEML2026}
Odrzyw\'{o}\l{}ek, A. (2026).
All elementary functions from a single binary operator.
\textit{arXiv preprint arXiv:2603.21852v2} [cs.SC].

\bibitem{hendrycksGELU2016}
Hendrycks, D., \& Gimpel, K. (2016).
Gaussian error linear units (GELUs).
\textit{arXiv preprint arXiv:1606.08415}.

\bibitem{browne2012survey}
Browne, C.~B., et al. (2012).
A survey of Monte Carlo tree search methods.
\textit{IEEE Transactions on Computational Intelligence and AI in Games},
4(1), 1--43.

\bibitem{sitzmann2020siren}
Sitzmann, V., Martel, J., Bergman, A., Lindell, D., \& Wetzstein, G.
(2020).
Implicit neural representations with periodic activation functions.
\textit{Advances in Neural Information Processing Systems},
33, 7462--7473.

\bibitem{paszke2019pytorch}
Paszke, A., et al. (2019).
PyTorch: An imperative style, high-performance deep learning library.
\textit{Advances in Neural Information Processing Systems},
32.

\end{thebibliography}

\end{document}
